\section{Background}
\label{Background}

\subsection{Test Gaming and Evaluation}

Scientific-agent benchmarks assess both discoveries and executable analyses.
DiscoveryBench evaluates hypotheses derived from datasets, while
ScienceAgentBench evaluates programs against scientific task
criteria~\cite{majumder2025discoverybench,chen2025scienceagentbench}.
ScienceAgentBench reports agents reading test labels instead of training the
requested model, showing that an expected answer can conceal an invalid
procedure~\cite{chen2025scienceagentbench}.

Weak tests also admit ordinary errors. PatchDiff finds behavioural differences
between agent and developer patches on SWE-bench Verified, with manual review
confirming errors in a subset~\cite{wang2026solved}. Passing tests therefore
provides incomplete evidence of correctness, while failing additional checks
does not alone establish gaming. ImpossibleBench isolates specification-violating
shortcuts by introducing conflicting tests. Hiding tests reduces successful
exploitation but also reduces performance on the original tasks; read-only
access prevents test edits without eliminating
special-casing~\cite{zhong2025impossiblebench}. Test access can therefore affect
both legitimate implementation and opportunities for exploitation.

\subsection{Hidden Oracles and Test Adequacy}

Preventing exploitation and detecting it are different evaluation problems.
The EvilGenie preprint compares held-out tests, test-edit detection, and LLM
judges. Its held-out tests miss some heuristic solutions and flag some
legitimate attempts that mishandle edge cases~\cite{gabor2026evilgenie}.
This qualifies the role of hidden tests: restricting access may discourage
shortcuts, but detector accuracy still depends on what the tests exercise.

Barr et al.'s oracle survey distinguishes generating inputs from deciding
whether outputs are correct~\cite{barr2015oracle}. Metamorphic testing checks
relations between executions, reducing reliance on known answers. Segura et al.
show that relation selection affects fault detection~\cite{segura2016survey}.
For example, constant outputs satisfy row-order invariance; a suitable relation
requiring an output change can expose that shortcut. The relevant issue is
which behaviours a relation rules out, rather than how many tests are hidden.

Property-based tests check specified properties over generated inputs.
Ravi and Coblenz use mutation testing to evaluate property-based tests in 40
Python projects, finding higher average detection per test than unit tests,
with variation across property categories~\cite{ravi2025pbt}. Their
observational results support evaluating test design, without guaranteeing
superiority on new tasks. Mutation scores measure detection of injected code
changes, so evidence about agent shortcuts must also come from evaluating
those shortcuts directly.

\subsection{Agent Self-Review and Disclosure}

Language-based assessments require separate validation. Turpin et al. find
that reasoning explanations can omit influences on model answers, although
their experiments do not evaluate code self-review~\cite{turpin2023unfaithful}.
More directly, a METR research report describes task-dependent responses when
agents are questioned after reward hacking: some acknowledge violating user
intent, while others defend their behaviour~\cite{vonarx2025rewardhacking}.
This report is not peer-reviewed, and its examples do not establish general
self-review accuracy. Together, these findings motivate comparing admissions
with implementation evidence and distinguishing later acknowledgement from
disclosure during execution or in the final answer.

\subsection{Motivation for This Study}

The reviewed work establishes risks in test-based evaluation, but does not
settle which checks expose shortcuts tied to scientific datasets or whether
agents disclose those shortcuts consistently. We examine these questions
across six adapted tasks, separating ordinary and deliberately conflicting
test conditions. Dataset-anchored metamorphic checks extend evaluation beyond
fixed expected answers, while planted implementations and mutation testing
assess detector strength. Comparing these checks with disclosure at different
stages connects software correctness with the reliability of the agent's report.
